% $Id: developers_manual.tex  $

\chapter{Developer's manual}
\label{developers_manual}
\minitoc

\section{Introduction}

There is an increasing interest in applying machine learning methods to detect predictive patterns in neuroimaging data. The accumulative data size in neuroimaging, on the other hand, starts reaching a big data scale and opening opportunities for more engineers and data scientists interested in mental health and neuroscience to apply their expertise to this field. Collaborations between neuroscience and machine learning communities show promises to discover novel and informative biomarkers for brain diseases. To meet the needs from both communities, we add new contents about PRoNTo in this v3.0 manual for users who are interested in understanding it at a deeper level. Therefore, this chapter will describe the main framework of PRoNTo and its backend core functions in order to provide a roadmap for developers.
\\

This chapter will begin with descriptions on the folder structure, which explains how the files and library folders are organised, and some naming conventions we use. Then we will explain the five functional modules: `Data \& Design', `Prepare feature set', `Model: Specify new' and `Model: Specify from', `Model: Run', and finally `Compute weights'. Specifically, we will provide details for

\bi
\item How the GUI and Batch interfaces are built.
\item What functions are called to arrange the user inputs to PRoNTo accepted structures.
\item How the data structures are passed to core functions carrying on computations and model estimations.
\item How cross-validations are implemented.
\item How machines from third-party libraries are called.
\item How weights are computed for visualization.
\ei

We will also provide detailed information on specific functions implementing each module for reference. We think that understanding the structure of PRoNTo shall be enough for developers to implement and interface their own functionalities, hence further details on  particular components (e.g., implementing a new data type) beyond the framework backbone are not included here. 


\section{PRoNTo folder structure}

PRoNTo contains several subfolders under the main PRoNTo folder. In this section, we will provide a brief description of each subfolder.

\begin{itemize}
\item The \textit{\_devUtils} and \textit{\_unitTests} folders are particular tools for testing the software. The \textit{atlas} folder contains two atlases: the AAL atlas with labels, and the Brodmann atlas without labels. If the developer wishes to add atlases, this is the recommended place for that. Atlases in NIfTI format can be created easily by using SPM or manually by the user. Please refer to the last point in `Features' in the Section \ref{ssec:feature_nifti_mat} of the user manual. If you use multiple kernel learning (MKL) and want to see ROI contributions in `Display weights', a .mat file storing ROI labels should be created alongside with the atlas.  This .mat file should comprise a cell array, and each cell should contain a label of a specific region. Please refer to point `ROI Label' in Section \ref{ssec:add_plots} in the user manual for more details on how to create and use this .mat file. Atlases in .mat format can be created by users manually.
 
\item The \textit{batch} folder contains the files for the Batch module, which includes both the configuration (with `cfg' in the function names) and job execution (with `run' in the function names) functions. For each module in PRoNTo, we have corresponding configuration and execution files. The configuration functions generate the trees defining the interfaces of the modules in Batch. The job execution functions pass user inputs to lower level functions in PRoNTo for calculations and estimations. GUI and Batch share the same lower functions for computations and estimations. This modular design eases the development of new machines and functions in each module. 

\item The \textit{machines} folder contains the different regression and classification algorithms. There are functions from several third-party libraries (e.g., libsvm, liblinear and SimpleMKL), `machines' which are wrapper functions (with `prt\_machine' in the function names) interfacing the libraries to PRoNTo, and functions computing weights for each machine.  If the developer wishes to add a new machine, they can open these functions and read their input and output arguments. There are a few naming conventions: `gp' – Gaussian processes, `cla' – classification, `krr' – kernel ridge regression. Names of these algorithms also refer to libraries, for example, `gp' refers to `gpml' library, `MKL' refers to simpleMKL library.

\item The \textit{manual} folder contains latex files for the manual, images and other contents. Developers can add to the documentation here.

\item The \textit{masks} folder contains one first-level mask in the Data \& Design module until before version 3. This changed in PRoNTo v3.0 and more masks (in NIfTI format) can be added.

\item In \textit{utils} folder, there are several lower level functions e.g. checking if inputs are alphanumeric, or performing operations on kernels, such as splitting data in the kernel and normalising kernels. 
 
\item Beyond the above mentioned subfolders, functions with `ui' in their names define the GUI interfaces, i.e. *\_ui\_*.fig and *\_ui\_*.m file pairs define the figure and how the GUI behaves, respectively. We used GUIDE (GUI development environment) to create GUI and add custom functions and codes to extend these interface functions. 
 
\item The rest of the `prt\_' functions perform the computations and estimations defined in each module.

\end{itemize}

%======================================================================



%======================================================================

\section{Data \& Design}

\subsection{PRT fields created}

The PRT structure is an important data structure (a .mat file) created and repeatedly used by PRoNTo. It is firstly created after the `Data \& Design' module is initiated / specified by the user. At first, it consists of two fields: `group' and `masks'.  Users input all their data, experimental design (e.g., modalities, runs, conditions) and other related information here (e.g. covariates, regression values). After this step, data information is stored in the group field of PRT.mat, and the first level masks,  HRF overlap and delay parameters are stored in the masks field of the PRT.mat. Below is a list of subfields in \textcolor{blue}{PRT.group} and \textcolor{blue}{PRT.masks}:

\begin{itemize}
\item \textcolor{blue}{PRT.group}

\begin{itemize}
\item \textcolor{blue}{PRT.group.gr\_name}: Storing group names.
\item \textcolor{blue}{PRT.group.subject}: Storing subject information.

\begin{itemize}
\item \textcolor{blue}{PRT.group.subject.subj\_name}: Storing subject names.
\item \textcolor{blue}{PRT.group.subject.modality}: Subfields containing information on regression targets, covariates, modality name, design and scans. Note: classification labels are specified in `Model: Specify new/from' modules later.


\end{itemize}
\end{itemize}

\item \textcolor{blue}{PRT.masks}

\begin{itemize}
\item \textcolor{blue}{PRT.masks.mod\_name}: Name of corresponding modality.
\item \textcolor{blue}{PRT.masks.type}: Storing the data type.
\item \textcolor{blue}{PRT.masks.fname}: The full path of the specified mask.
\item \textcolor{blue}{PRT.masks.hrfoverlap}: Storing information on HRF overlap (default: 0).
\item \textcolor{blue}{PRT.masks.hrfdelay}: Storing information on HRF delay (default: 0).

\end{itemize}
\end{itemize}


%----------------------------------------------------------------------

\subsection{Files created}

PRT structure, saved in a PRT.mat file in the path specified by the user.


%----------------------------------------------------------------------

\subsection{GUI behaviour}

GUI of Data \& Design module provides interfaces for users to specify where to store the PRT structure and files, data information, experimental design, masks and all other related information.

\begin{itemize}
\item Save button creates the PRT structure.
\item Load button loads previously created PRT.mat files and checks their compatibility.
\item Review button provides a graphical review of all the input information and the interface for modifying some HRF parameters.
\item \textcolor{red}{prt\_ui\_main.m} and \textcolor{red}{prt\_ui\_main.fig} function pair define the GUI of the main window.
\item Button click on `Data \& Design' in the main window will call \textcolor{red}{prt\_ui\_design.m} and \textcolor{red}{prt\_ui\_design.fig} to open the Data \& Design module.
\end{itemize}



%----------------------------------------------------------------------

\subsection{Batch behaviour}

The Batch of Data \& Design module provides users similar functions and interfaces as GUI. However, there are three major differences.

\begin{itemize}
    \item Firstly, the batch job can be saved as a .mat file or a .m function. Both users and developers can make changes directly in these files.
    \item Secondly, both users and developers should be careful about the names of modalities when using the Batch, i.e. the names needs to be consistent.
    \item Finally, the HRF delay and overlap can be changed directly in the Batch of Data \& Design. In GUI, they can be modified in Review.
\end{itemize}

%----------------------------------------------------------------------

\subsection{Functions called}

The GUI functions related to the Data \& Design module can be found under the root PRoNTo folder, whereas the Batch functions of this module are under the \textit{batch} folder.
\\

\noindent
\textbf{\underline{GUI:}}

\begin{enumerate}

\item \textcolor{red}{prt\_ui\_design.m} and \textcolor{red}{prt\_ui\_design.fig}

These functions define GUI interface and other functions for the Data \& Design module. The function `\textcolor{red}{prt\_ui\_design.m}' initializes the GUI window, executes the commands associated with each button, checks if the inputs are correct, and creates outputs in \matlab{} prompt. If the developer chooses to make changes in this GUI module (e.g., a button’s or menu’s functions), you can find related sub-functions with corresponding names in \textcolor{red}{prt\_ui\_design.m}. Some important functions called by \textcolor{red}{prt\_ui\_design.m} are summarized as follows.


\begin{itemize}
\item \textcolor{red}{prt\_data\_modality.m} and \textcolor{red}{prt\_data\_modality.fig}

The two files define the GUI panel for specifying the information of the modality, experimental design, covariates and regression targets. The \textcolor{red}{prt\_data\_modality.m}  also calls:

\begin{itemize}
\item \textcolor{red}{prt\_data\_conditions.m} and \textcolor{red}{prt\_data\_conditions.fig}, for receiving inputs for conditions, i.e. experimental design.
\item \textcolor{red}{prt\_check\_design.m}, for checking the design and discards scans according to overlapping conditions or the minimum time interval between conditions based on the HRF function and specified HRF parameters.
\item \textcolor{red}{prt\_data\_targets.m} and \textcolor{red}{prt\_data\_targets.fig}, for opening the target window to examine and/or modify targets.
\item \textcolor{red}{prt\_get\_design\_MEEG.m}, for filling the design from MEEG file. In v3.0, new data types (.mat and MEEG) are allowed, hence relevant codes are added to interface functions, and/or core functions carrying out the backend computations. 

\end{itemize}

\item \textcolor{red}{prt\_get\_defaults.m} and \textcolor{red}{prt\_defaults.m}

Default settings of PRoNTo are stored in \textcolor{red}{prt\_defaults.m}, including global defaults such as colour, parameters for the data and design, preprocessing, default atlas for ROI definition, arguments of different machines, and parallelization of the code. The function \textcolor{red}{prt\_get\_defaults.m} gets the relevant default values from \textcolor{red}{prt\_defaults.m} for this module. For instance, the code `color = \textcolor{red}{prt\_get\_defaults}(`color')' sets colors of backgrounds and figure parameters for Data \& Design module. \textcolor{red}{prt\_defaults.m} is constantly called by many functions in PRoNTo.

\item \textcolor{red}{prt\_data\_review.m} and \textcolor{red}{prt\_data\_review.fig}

These two functions define the GUI for reviewing the data.

\end{itemize}

\item \textcolor{red}{prt\_load.m} and \textcolor{red}{prt\_struct.m}

This function is called when clicking the Load button. It loads PRT.mat and checks its fields for backwards compatibility. If you wish to add a field to the PRT structure, and that field is necessary for display or further analyses, the \textit{prt\_struct} file should be amended to set default values for that field to ensure that older PRT files can still be read and modified. 

\end{enumerate}

\noindent \textbf{\underline{Batch:}}

\begin{enumerate}
    \setcounter{enumi}{2}

\item \textcolor{red}{prt\_cfg\_design.m} and \textcolor{red}{prt\_run\_design.m}

\begin{itemize}
\item The function \textcolor{red}{prt\_cfg\_design.m} is the Data \& Design module configuration file that builds the PRT.mat data and design structure.
\item The function \textcolor{red}{prt\_run\_design.m} is the job execution function, which takes a harvested job data, rearranges them into PRT structure and saves it.
\end{itemize}

\end{enumerate}

At this stage, the GUI and batch have very little code in common. Hence, changes in the GUI have to be performed independently in the batch and conversely.  


%======================================================================



%======================================================================

\section{Prepare feature set}

\subsection{PRT fields created}

There are two fields created in PRT.mat in this step, too. Field \textit{fs} (feature set) and \textit{fas} (file array structure). The information for the two fields is specified by users in the `Specify modality to include' panel.
\\


For each selected modality, a file array is created containing the data (detrended if specified). The field \textcolor{blue}{PRT.fas} is linked to this file array. In \textit{fas} field, the information stored refers to the modality (as defined in Data \& Design). This means that multiple MRI sessions, entered as different runs, will each create a specific \textit{fas}. A \textit{fas} is linked to a \textit{.dat} file created at this stage, which is with dimensions of \#images x \#features within the first-level mask. A \textit{fas} structure also stores feature operations (detrend and detrend parameters), header of the first image for recovering images if needed, and indices of all features within the first-level mask. \\

The field \textit{fs} refers to the kernel/feature set that is created. It is a structure of size number of feature sets created by the user. This structure of \textit{fs} stores information on the feature set name, kernel file, ID matrix, information of using multiple kernels, second-level mask and atlas. This field also contains information about the second-level mask and/or the atlas (used to build the kernel).
\\

One important subfield is the ID matrix, which contains information about the modalities selected in the feature set arranged in six different columns: group, subject, modality, condition, block and scan.

\begin{itemize}
\item The name of the columns are stored in \textcolor{blue}{PRT.fs.id\_col\_names}.
\item The group numbers, subject numbers, modality labels, condition numbers, block labels and scan numbers themselves are stored in \textcolor{blue}{PRT.fs.id\_mat}, whose dimension is the number of samples x 6. Columns in \textcolor{blue}{PRT.fs.id\_col\_names} and \textcolor{blue}{PRT.fs.id\_mat} are in corresponding order.
\end{itemize}

The ID matrix is referred to in the `Review kernel \& CV' reviewing tool in PRoNTo, where a graphical display shows the structure of this matrix.

\begin{itemize}
\item In addition, there is a \textit{fas} subfield in \textcolor{blue}{PRT.fs}, that is different from \textcolor{blue}{PRT.fas}. This subfield in \textit{fs} refers to the modality and scan indices of the feature set modalities in the \textit{fas} structure. It is useful when modalities are concatenated in samples (e.g. accessing the correct scans from multiple runs in an MRI experiment, with each run linked to a different \textcolor{blue}{PRT.fas} file array).

\end{itemize} 

As mentioned in the Data \& Design module, in v3.0, different data formats (.mat, NIfTI and MEEG) are supported. Users can specify which data type and modality shall be used to build the feature sets. Some changes according to these are made in both GUI and Batch of this module, which will be more described as follows. 



% \begin{itemize}
% \item \textcolor{blue}{PRT.fs}

% \begin{itemize}
% \item \textcolor{blue}{PRT.fs.fs\_name} - 
% \item \textcolor{blue}{PRT.fs.k\_file} - 
% \item \textcolor{blue}{PRT.fs.id\_col\_name} - 
% \item \textcolor{blue}{PRT.fs.} - 
% \item \textcolor{blue}{PRT.fs} - 
% \item \textcolor{blue}{PRT.fs} - 
% \item \textcolor{blue}{PRT.fs} - 


% \end{itemize}

% \item \textcolor{blue}{PRT.masks}

% \end{itemize}


%----------------------------------------------------------------------

\subsection{Files created}

\begin{itemize}
\item FAS .dat file(s).
\item A kernel matrix (.mat with variable Phi, a cell array with a kernel in each cell).
\item Updated masks and/or atlases (NIfTI resampled at the size of the input data images).
\end{itemize}


% In this module the files created are the following:

% \bi
% \item .dat file(s)
% \item A kernel matrix (.mat with variable Phi, a cell array with a kernel in each cell)
% \item Updated masks
% \item And/or atlases (NIfTI resampled at the size of the input data images)
% \ei

The .dat file refers to one \textcolor{blue}{PRT.fas} structure and stores all features in the first-level mask for each modality. It is a binary file. If computing multiple kernels using a second-level atlas is chosen by the user, the kernel matrix will contain multiple cells. Each cell stores information of one kernel. 

%----------------------------------------------------------------------

\subsection{GUI behaviour}

The GUI of this module takes the PRT.mat created in the Data \& Design module. It provides interfaces for users to further specify feature sets, their names, selected modalities, operations on the feature sets and atlas for defining kernels.
\\

In v3.0, users may specify three different data types (NIfTI, .mat and MEEG) in the Data \& Design module. If only the NIfTI or .mat file is specified, after loading PRT.mat, GUI calls for the ‘Specify modality to include’ panel. If more than two of the data types are specified, e.g. NIfTI and .mat, or all of them, GUI calls for ‘Specify type of data’ after loading PRT.mat. The reason for this design is that ‘Specify modality to include’ panel is different between MEEG and the other two types. The function \textcolor{red}{prt\_ui\_main.m} checks the type of data (modalities) in the loaded PRT.mat and calls for these panels. Buttons on the panels have names of the modalities. Their functions are defined using corresponding GUI files, such as \textcolor{red}{prt\_ui\_fs\_alldatatype.m} and \textcolor{red}{prt\_ui\_fs\_alldatatype.fig}. Sub-functions like \textcolor{red}{prt\_ui\_fs\_alldatatype.m} will gather modality information based on the button name, and pass it to \textcolor{red}{prt\_ui\_prepare\_data.m}. This function then calls for the right ‘Specify modality to include’ panel for the selected data type/modality. This panel accepts user inputs on which modalities to be included and several other related inputs, then ‘Prepare feature set’ window is called. After filling in the fields of this window by users, \textcolor{red}{prt\_ui\_prepare\_data.m} passes user inputs to \textcolor{red}{prt\_fs.m} function.
\\

Main functions in the Prepare feature set module, such as building the file array and computing feature set matrices (including multiple kernels) are completed by \textcolor{red}{prt\_fs.m} or \textcolor{red}{prt\_fs\_EEG.m}. Here, we take \textcolor{red}{prt\_fs.m} as an example for some further illustrations. This function calls \textcolor{red}{prt\_init\_fs.m} to initialize the file arrays, kernels and feature set parameters. Then it calls \textcolor{red}{prt\_fs\_modality.m} to build file arrays and kernels. The .dat files linked to file arrays are created here.  A sub-function within \textcolor{red}{prt\_fs.m} named as \textcolor{red}{prt\_compute\_ROI\_kernels} builds kernels based on the second-level atlas. 


%----------------------------------------------------------------------

\subsection{Batch behaviour}

The `Prepare feature set' module in the Batch is similar to GUI. However, the batch requires consistent names across modules. Running important modules altogether will overwrite the PRT.mat, hence delete the links between the PRT.mat and the computed feature set(s) and kernel(s). \textit{FAS} will be recreated each time the Batch is launched.  Function \textcolor{red}{prt\_cfg\_fs.m} defines the tree of this Batch module. The function \textcolor{red}{prt\_run\_fs.m} harvests user inputs from the tree. It checks the modality and data types, then calls \textcolor{red}{prt\_fs.m} or \textcolor{red}{prt\_fs\_EEG.m} to execute computations on file arrays, kernels and feature set. After loading PRT.mat, types of data can be selected under ‘Data format’ field. 


%----------------------------------------------------------------------

\subsection{Functions called}

\textbf{\underline{Core:}}

\begin{enumerate}

\item \textcolor{red}{prt\_fs.m}

The function gathers information on the modalities to build file arrays containing the (detrended) data and computes the linear kernels. It resizes the 2nd level masks and atlases if needed. The 2nd-level masks are used here as an important tool. It is used to define more precise scope for building the feature sets, for example, restricting the analysis to certain brain regions. Multiple kernels can be built for multiple modalities, or multiple regions of a specified atlas. \textcolor{red}{prt\_fs.m} calls  \textcolor{red}{prt\_fs\_modality.m} directly to compute multiple kernels for different modalities, but calls \textcolor{red}{prt\_compute\_ROI\_kernels} to build ROI defined kernels for multiple regions in the atlas.  A list of functions \textcolor{red}{prt\_fs.m} calls is as follows:

\begin{itemize}
\item \textcolor{red}{prt\_init\_fs.m} to populate basic fields in \textcolor{blue}{prt.fs}, like the file array structure, kernel data structure and feature set parameters. 
\item \textcolor{red}{prt\_fs\_modality.m} to write the kernel matrix to the user-defined path as well as the feature set if needed.
\item \textcolor{red}{prt\_compute\_ROI\_kernels}. This is a sub-function within \textcolor{red}{prt\_fs.m}. It reads voxel indexes within both the 2nd-level mask and the 1st-level mask first. It then computes kernel for each region in the user-specified atlas and stores indices in the image for the weight computation.   
\end{itemize}

\item \textcolor{red}{prt\_init\_fs.m}

The function initialises the file arrays, kernels and feature set parameters.
     
\item \textcolor{red}{prt\_fs\_modality.m}

Code in this function are divided in two parts: building file arrays and computing kernels. If a file array has been already built for a given modality, it won’t be built again, but read from the existing file. Operations like scaling are applied on the kernel only. To build multiple kernels, \textcolor{red}{prt\_fs.m} calls \textcolor{red}{prt\_fs\_modality.m} as many times as the number of kernels defined. It calls \textcolor{red}{prt\_get\_defaults.m} to get defaults of \textit{fs}, and \textcolor{red}{prt\_load\_blocks.m} to access one or more blocks of data to avoid memory overload. 

\end{enumerate}


\noindent \textbf{\underline{GUI:}}

\begin{enumerate}
\setcounter{enumi}{3}

\item \textcolor{red}{prt\_ui\_prepare\_data.fig} and \textcolor{red}{prt\_ui\_prepare\_data.m}

This pair of functions define the GUI to specify how many feature sets will be used. These modalities can either be concatenated (e.g. multiple runs of an experiment), or combined in multiple kernel settings (e.g. sMRI grey and white matter). For each modality entered, the function calls \textcolor{red}{prt\_ui\_prepare\_datamod.m} to let users specify parameters for the selected modality, or data types in .mat or NIfTI formats. It calls \textcolor{red}{prt\_ui\_prepare\_dataMEEG.m} and \textcolor{red}{prt\_ui\_prepare\_dataMEEG.fig} to let users specify parameters for MEEG data type.

\item \textcolor{red}{prt\_ui\_prepare\_datamod.fig} and \textcolor{red}{prt\_ui\_prepare\_datamod.m}
 
This pair of files are called by \textcolor{red}{prt\_ui\_prepare\_data.m} to accept user inputs about modalities, NIfTI and .mat data types to compute the file array, kernels and feature set. 

\item \textcolor{red}{prt\_ui\_prepare\_dataMEEG.m} and \textcolor{red}{prt\_ui\_prepare\_dataMEEG.fig}

This pair of files are called by \textcolor{red}{prt\_ui\_prepare\_data.m} to accept user inputs about MEEG data to compute file array, kernels and feature sets. 


\end{enumerate}

\noindent \textbf{\underline{Batch:}}

\begin{enumerate}
\setcounter{enumi}{6}

\item \textcolor{red}{prt\_cfg\_fs.m} and \textcolor{red}{prt\_run\_fs.m}

The \textit{fs} configuration file constructs the tree structure for this module. It provides user interfaces for specifications of the feature set for each data type and modality. The job execution function takes all inputs from \textcolor{red}{prt\_cfg\_fs.m}. It arranges them into proper data structure and stores this structure in a variable called `input'. It then passes PRT.mat and `input' to \textcolor{red}{prt\_fs.m} or \textcolor{red}{prt\_fs\_EEG.m} to build file arrays and compute kernels.


\end{enumerate}




%======================================================================



%======================================================================

\section{Model: Specify new/from}

In v3.0, the `Specify model' module of versions 2.X is split to `Model: Specify new' and `Model: Specify from' modules. They compute the inputs for the models but don't estimate/run the models. The main inputs for this module are: feature sets, where users select which feature set will be used to build the model, whether to use kernel or not, model type (regression or classification), selection of samples and classes, machines, cross validation (CV) strategy and data operations. `Model: Specify new' is used when users want to specify a new model to run. `Model: Specify from' is used when users want to specify models similar to previous ones.

\subsection{PRT fields created}

\textbf{\underline{Model: Specify new}}
\\

Two subfields in \textcolor{blue}{PRT.model} are created and filled: \textcolor{blue}{PRT.model.model\_name}, which stores the name of the model and \textcolor{blue}{PRT.model.input} which stores all the other information. Another subfield \textcolor{blue}{PRT.model.output} is only created but is not filled here. It will be filled after the model is estimated (`Model: Run' module). Below is an illustration of the information stored in \textcolor{blue}{PRT.model.input}:
 
\begin{itemize}
 
\item \textcolor{blue}{PRT.model.input.use\_kernel}: Whether to use kernel methods (1) or not (0, for non-kernel methods). Non-kernel methods are available from v3.0.

\item \textcolor{blue}{PRT.model.input.type}: Regression or classification. This field is called very often within PRoNTo.

\item \textcolor{blue}{PRT.model.input.machine}: Stores information of the chosen machine. The algorithms can be found in the subfolder \textit{PRoNTo/machines/}. More details on how to call the machines can be found in the following `Model: Run' section. 

\item \textcolor{blue}{PRT.model.input.use\_nested\_cv}: Whether to optimize hyperparameters (1) or not (0).

\item \textcolor{blue}{PRT.model.input.nested\_param}: A vector of user-defined parameters if \textcolor{blue}{PRT.model.input.use\_nested\_cv} is 1. In versions 2.X, this can also be empty, which means to leave it for the algorithm to set values by default. The default values are specified in \textcolor{red}{prt\_nested\_CV.m} at the `Model: Run' step. In v3.0, this has changed. All default parameters are automatically filled in the corresponding fields of GUI or Batch, once `Optimize hyperparameter' field is activated. Hence from v3.0 they are taken as inputs by default from user interfaces.

\item \textcolor{blue}{PRT.model.input.cv\_type\_nested}: The CV strategy for the inner CV loop.

\item \textcolor{blue}{PRT.model.input.cv\_k\_nested}: Value of k for the k-fold inner CV loop. For instance, if 10 is input, it means that the strategy takes 90\% of the data for training, and 10\% of the data for test.
\begin{itemize}
\item Leave-one-out has k=0.
\end{itemize}

\item \textcolor{blue}{PRT.model.input.subsample}: Whether to subsample the different classes accordingly or not, in order to have balanced classes.

% \item \textcolor{blue}{PRT.model.input.indmodels}:

\item \textcolor{blue}{PRT.model.input.class}: Contains the user-specified selection for each class. 

\begin{itemize}
\item It's a structure of size number of classes (classification) or 1 (regression). It specifies the groups, subjects within that group and conditions for each class, with a similar structure as in \textcolor{blue}{PRT.group} (without the modalities). And we identify which indices in the global ID matrix are selected to build the specified class.
\item The CV matrix and everything else are built from here. This can be viewed in Review kernel \& CV. Note: it’s always good to have a check on the CV matrix to ensure the inputs do what you want. You can also make a CV matrix yourself and input via Custom choice in Cross-validation scheme. In the Review module, you can view the classes specified.
\end{itemize}


\item \textcolor{blue}{PRT.model.input.samp\_idx}: Index of the selected samples for this model, in the corresponding feature set.

% \item \textcolor{blue}{PRT.model.input.include\_allscans}:


\item \textcolor{blue}{PRT.model.input.targets}: Prediction targets for building the model.

% \item \textcolor{blue}{PRT.model.input.targ\_allscans}:

\item \textcolor{blue}{PRT.model.input.covar}: Values of the confounders to be regressed out.

% \item \textcolor{blue}{PRT.model.input.cov\_allscans}:

\item \textcolor{blue}{PRT.model.input.cv\_k}: Value of k for the k-fold outer CV loop (i.e. the main or external CV).

\item \textcolor{blue}{PRT.model.input.cv\_mat}: CV matrix for the external CV.

\item \textcolor{blue}{PRT.model.input.cv\_type}: CV type for the external CV.

\item \textcolor{blue}{PRT.model.input.operations}: Stores user selected data/kernel operations as an integer vector. The operations will be applied according to this order. 
 
\end{itemize} 

\noindent \textbf{\underline{Model: Specify from}}
\\

`Model: Specify from' module allows users to define models similar to previously specified models,  instead of filling in all parameters multiple times. In GUI, this module is implemented in \textcolor{red}{prt\_ui\_copy\_model.m} and \textcolor{red}{prt\_ui\_copy\_model.fig}. In Batch, fields that can be modified are the feature sets, model type and data operations. The batch functions define this module are \textcolor{red}{prt\_cfg\_copy\_model.m} and \textcolor{red}{prt\_run\_copy\_model.m}. `Model: Specify from' creates another structure in \textcolor{blue}{PRT.model}. Similar to `Model: Specify new', the model name and input fields are filled after this step. All the subfields inside \textcolor{blue}{PRT.model}.input structure are the same to those created by `Model: Specify new'.
\\

The two modules create and estimate inputs for estimating the model. The higher-level users and developers can copy \textcolor{blue}{PRT.model} and change any of the sub-fields listed above when it is needed manually. CV matrix is computed based on CV schemes specified by the user.

%----------------------------------------------------------------------

\subsection{Files created}

No new files are created at this stage.


%----------------------------------------------------------------------

\subsection{GUI behaviour}

\textbf{\underline{Model: Specify new}}
\\

In GUI, users can choose only to specify a model, or to specify and run a model. It computes the sample indices based on the classes, computes CV matrix based on the ID matrix (and labels if `Leave-per-Group-Out') and CV specifications, and it gathers the targets and covariates based on the class definition. ID matrix is used here as a reference to obtain these sub-groups of indices, targets and covariates for building a particular model. CV matrix can also be input by users via Custom CV. In v3.0, it also contains default hyperparameters of the selected machines for building the model. The function pair of \textcolor{red}{prt\_ui\_model.m} and \textcolor{red}{prt\_ui\_model.fig} define the GUI, arrange user inputs into appropriate structures and send them to several core functions (e.g., \textcolor{red}{prt\_model.m}) to implement the above computations. \\

\noindent \textbf{\underline{Model: Specify from}}
\\

In GUI, it is the same as the Specify new module, users can choose Specify model alone or Specify and run model. After loading PRT.mat, users can select  which previously built model is used as a reference. Some parts of the window is automatically filled, i.e. these information is copied from the selected referential models. For example, model type, and the main CV scheme. Users can change other activated fields, such as feature sets, machines, inner CV scheme, and data operations. Since this module copies information from the referential model, it does not carry out computations described in Specify new module. User defined changes for the model will be filled in corresponding fields in \textcolor{blue}{PRT.model}. The function pair of \textcolor{red}{prt\_ui\_copy\_model.m} and \textcolor{red}{prt\_ui\_copy\_model.fig} define this module, organise user inputs, and updates PRT. 

%----------------------------------------------------------------------

\subsection{Batch behaviour}

\textbf{\underline{Model: Specify new}}
\\

Batch provides similar interfaces and functions to the GUI, but the Specify model  and run model modules are separated. This module calls \textcolor{red}{prt\_model.m} to configure and build \textcolor{blue}{PRT.model} structure too. \textcolor{red}{prt\_cfg\_model.m} and \textcolor{red}{prt\_run\_model.m} carry out this module in batch. 
\\

\noindent \textbf{\underline{Model: Specify from}}
\\

Unlike GUI, the list of built models is not automatically shown in the batch. Users shall input the exact name of the previous model to copy from. Three fields can be changed: Feature sets, Model type and Data operations. This module calls \textcolor{red}{prt\_run\_copy\_model.m} to update \textcolor{blue}{PRT.model} structure.

%----------------------------------------------------------------------

\subsection{Functions called}

\textbf{\underline{Core:}}

\begin{enumerate}

\item \textcolor{red}{prt\_model.m}

This function initialises the model data structure (by calling \textcolor{red}{prt\_init\_model.m}), computes targets and the sample index.

\begin{itemize}
\item In PRoNTo v2.1, classification and regression are addressed differently. In classification, the codes visit information stored in the design. For regression, only subject-level analysis is available, so the codes do not go into the design.
The sub-functions \textit{compute\_targets} and \textit{compute\_target\_reg} inside \textcolor{red}{prt\_model.m} compute the prediction targets for classification and regression, respectively. Targets that are not selected for this model are put to zeros. From PRoNTo v3.0 both classification and regression compute the prediction targets using an updated version of the \textit{compute\_targets} sub-function.
\end{itemize}
 
One note is that the labels of different classes are 1,2,3,..., here referring to Class 1, Class 2, Class 3, ... respectively. They will be converted to class labels accepted by the actual machines inside the machine function, if needed. For example, if Faces and Houses are selected by the user as Class 1 and Class 2, respectively, and if we use SVM, the labels 1 and 2 will be converted to -1 and 1 in the machine \textcolor{red}{prt\_machine\_svm\_bin.m}.

\item \textcolor{red}{prt\_compute\_cv\_mat.m}

It computes the CV matrix. Inner and external CV loops call the same function. This function is accessed in the Specify model and `Model: Run' modules. In PRoNTo, CV schemes are shown in abbreviations, for example:

\begin{itemize}
\item LOSO: leave-one-subject-out. If we have 56 subjects, and 10 folds. Then 5 subjects in each fold, the remaining 6 will be distributed to the first 6 folds. As a result, 6 subjects in the 1 to 6 folds, 5 subjects in the remaining 4 folds. The reason to do this is for distributing subjects more evenly.
\item LORO:leave-one-run-out. There is no leave-multiple-run-out option in PRoNTo, only leave-one-run-out.
\end{itemize}

Other CV schemes can be found in Section \ref{sec:cross_val} in this manual. PRoNTo also provides interfaces for users to input their own CV design, i.e. the custom CV. Users can input their own CV matrix via this interface. The custom CV saves the model and calls the \textcolor{red}{compute\_cv\_mat.m} if the user chooses a `basis'. Please note that nested CV will not be computed in `Model: Specify new' step, but during the model estimation.
 
\end{enumerate}


\noindent \textbf{\underline{GUI:}}

\begin{enumerate}
\setcounter{enumi}{2}

\item \textcolor{red}{prt\_ui\_model.m}  and \textcolor{red}{prt\_ui\_model.fig}

The GUI functions that define the interfaces of Specify model module. It calls several other functions among which:


\item \textcolor{red}{prt\_ui\_select\_class.m} and \textcolor{red}{prt\_ui\_select\_class.fig}

This function provide interfaces to users to select classes for classification problems.


\item \textcolor{red}{prt\_ui\_select\_reg\_new.m} and \textcolor{red}{prt\_ui\_select\_reg\_new.fig}

This function provide interfaces to users to select data for regression analysis in regression problems.

\item \textcolor{red}{prt\_ui\_specify\_CV\_basis.m} and \textcolor{red}{prt\_ui\_specify\_CV\_basis.fig}

This function provides interfaces for users to input their custom CV information. It lets users load their own CV design, or build their own CV based on several basic CV schemes available in PRoNTo, such as LOBO. These inputs serve as   `basis' for users to further modify CV design manually in a pop-up window after this step where \textcolor{red}{prt\_ui\_specify\_CV\_basis.m} calls \textcolor{red}{prt\_ui\_custom\_CV.m} and \textcolor{red}{prt\_ui\_custom\_CV.fig} to allow more custom CV specification manually.

\item \textcolor{red}{prt\_model.m}

This function configures and builds the \textcolor{blue}{PRT.model} structure.

\item \textcolor{red}{prt\_cv\_model.m}

This function runs a CV on a given model.

\item \textcolor{red}{prt\_ui\_copy\_model.m} and \textcolor{red}{prt\_ui\_copy\_model.fig}

This function pair defines the `Model: Specify from' module. 





\end{enumerate}


\noindent \textbf{\underline{Batch:}}


\begin{enumerate}
\setcounter{enumi}{9}

\item \textcolor{red}{prt\_cfg\_model.m} and \textcolor{red}{prt\_run\_model.m}

The configuration function defines the tree of the `Model: Specify new' module and the interface behaviours in Batch. It provides interfaces for users to specify which feature sets to use, the model type (classification and regression), machine type (kernel or non-kernel), inner and outer CV schemes, and data operations. Default values for machine hyperparameters and arguments are automatically filled in correspondent fields. Users can change these values. It gathers these information, then passes them to \textcolor{red}{prt\_run\_model.m}. \textcolor{red}{prt\_run\_model.m} takes the user inputs and arranges them into a proper model structure. It passes PRT.mat and this model structure to \textcolor{red}{prt\_model.m}. 

\item \textcolor{red}{prt\_cfg\_copy\_model.m} and \textcolor{red}{prt\_run\_copy\_model.m}
The configuration file defines the tree of the Specify from module and the interface behaviours in Batch. It provides interfaces for users to choose which previously built model to copy, and which fields in the model to modify. It passes these information to \textcolor{red}{prt\_run\_copy\_model.m}, which arranges these user inputs to a model structure with the new model name (user-specified). \textcolor{red}{prt\_run\_copy\_model.m} stores this new model to and updates PRT.mat. 


\end{enumerate}




%======================================================================



%======================================================================

\section{Model: Run}

`Model: Specify new' and `Model: Specify from' modules set up the configuration of the model and build the model structure. `Model: Run' module performs the training and testing. It includes the selection of which model to run, hyperparameter optimization by using inner CV schemes, model training and testing using external CV schemes, computing statistics for model performance measures, and PRT.mat updates. In v3.0, we have moved permutations from the `Display results' to this module. In GUI, this module can be run directly from specify model stage. In Batch, it is separated from the `Model: Specify new/from' modules. 
\\ %[0.3cm]

\noindent \textbf{\underline{Cross-validation}}
\\ %[0.15cm]

CV schemes are important to model training and testing. In PRoNTo, this is mainly performed by \textcolor{red}{prt\_cv\_model.m}, which executes both external and inner CV. After some initialization steps, it begins the main/external CV loops. If the user chooses to optimize hyperparameters, the external CV sends data information of each fold to the inner CV. Users can choose different CV strategies for inner and external CV. The function \textcolor{red}{prt\_nested\_cv.m} is called to optimize hyperparameters using the inner CV scheme. After this, the best hyperparameter is returned to \textcolor{red}{prt\_cv\_model.m} within each external fold. Whether the hyperparameter is optimized or not, one model is estimated per fold by calling \textcolor{red}{prt\_cv\_fold.m}. \textcolor{red}{prt\_cv\_model.m} passes PRT.mat file and a data structure \textit{fdata} to \textcolor{red}{prt\_cv\_fold.m}. The \textit{fdata} structure contains information within each fold on ID matrix, model index, CV matrix, class labels or targets, and kernels. \textcolor{red}{prt\_cv\_fold.m} returns the estimated model, related parameter values and targets. This is followed by statistic computations for model performance measures per external fold, by calling \textcolor{red}{prt\_stats.m}. In \textcolor{red}{prt\_cv\_model.m}, statistics are calculated and stored at fold level and model level. Model level statistical metrics are obtained by concatenating results across folds in versions 2.X, but by averaging metrics across folds in v3.0. PRT.mat is updated and saved at the end of \textcolor{red}{prt\_cv\_model.m}.
\\

\noindent \textbf{\underline{Hyperparameter optimization}}
\\

In many situations, it is preferable to optimize hyperparameters to increase the predictive model performance. If users choose this step, each external fold constructs ‘fdata’ structure and sends it alongside with PRT.mat to \textcolor{red}{prt\_nested\_cv.m}, which implements the inner CV loop. The training data in the external fold is splitted into training and testing sets inside \textcolor{red}{prt\_nested\_cv.m} according to the user-specified inner CV scheme. Before looping over inner CV folds, \textcolor{red}{prt\_nested\_cv.m} checks hyperpameter ranges. Hence, this function first loops over hyperparameters within the range and then loops over inner folds for each hyperparameter. It calls \textcolor{red}{prt\_cv\_fold.m} to perform the model training and testing. One model is created for each inner fold, generating predictions. Within every loop using a certain hyperparameter, model performance metrics are calculated for each inner fold. Statistic calculation is done by calling \textcolor{red}{prt\_stats.m}. In version 2, we concatenate predictions from each inner fold to obtain the model-level performance measures. In version 3, we average metrics across folds. Among the metrics, we use balanced accuracy for classification and MSE for regression. This process generates one metric value per hyperparameter. After completing the inner CV for all values of the hyperparameter, we choose the best one according to the metrics, i.e. the hyperparameter generating the maximum balanced accuracy is selected as the best one for classification and the hyperparameter generating the minimum MSE is selected as the best one for regression. \textcolor{red}{prt\_nested\_cv.m} returns the best hyperparameter to \textcolor{red}{prt\_cv\_model.m} for the outer model training and testing. In v3.0,  hyperparameter optimization can be performed on models with more than one hyperparameters.   
\\

\noindent \textbf{\underline{The machines}}
\\

Model training and testing for both external and inner CV are done by calling \textcolor{red}{prt\_cv\_fold.m}. This function configures model parameters and applies user-specified data operations before feeding them to the machines (e.g., \textcolor{red}{prt\_prepare\_task\_input\_STL.m}). It calls \textcolor{red}{prt\_machine.m} which will then distribute to the specific machine after some input checks. The machine-specific functions are in the \textit{machines} folder, with ‘machine’ in their names. Some of them are wrapper functions of the core functions in several third-party libraries, such as Libsvm, Liblinear and SimpleMKL. PRoNTo uses wrapper functions to organise inputs (data, class labels, targets, function arguments and hyperparameters) to formats accepted by actual machines in those libraries, and organize outputs from these machines to PRoNTo accepted formats. For instance, \textcolor{red}{prt\_machine\_svm\_bin.m} is a wrapper of binary SVM in Libsvm. PRoNTo takes the hyperparameter range from GUI or Batch, then combines it with other machine arguments required by Libsvm. Training and testing datasets are organized according to Libsvm conventions. Predictions and model parameters from the binary SVM machine are stored according to PRoNTo conventions. Machine arguments include options provided by third-party libraries and hyperparameters (e.g., C for SVM). We define machine options in \textcolor{red}{prt\_defaults.m}.  
\\

 
% \noindent \textbf{\underline{What are the input and output for a machine?}}
% \\

Generally speaking, each \textcolor{red}{prt\_machine\_*.m} function accepts two inputs, a structure d containing the data information and a structure args containing machine-specific information including options and hyperparameter values. In structure d: d.train and d.test contain the training and testing data respectively. The \textit{d.tr\_targets} refers to targets of the training data. There is no need to pass \textit{d.te\_targets} as model performance is evaluated outside of the machine, in PRoNTo \textcolor{red}{prt\_stats.m}. One flag d.use\_kernel contains information about whether the machine is a kernel machine or not. While this is functional from v3.0, in versions 2.X, this is by default always true. The output of each machine contains prediction values, function values from the decision function, model parameters and the type of the machine (classifier or regression). 
\\

\noindent \textbf{\underline{Permutations}}
\\

In v3.0, permutations have been moved to the `Model: Run' module. \textcolor{red}{prt\_permutation.m} implements the permutation test and saves the results. It accepts PRT.mat, number of permutations specified by the user, model ID (for which model to permute), path to save the results and whether to compute weights for models created by the permutations. If this option is selected, an extra field will be created inside \textcolor{blue}{PRT.model}, \textcolor{blue}{PRT.model.output.permutation}. This structure contains relevant information of models using permuted data, such as performance metrics, statistics, fold-level information (e.g., model parameters, function values, predictions) and the permuted data used for each permutation. Saving the permutation parameters is only optional. However if a user wants to build weights on the permutations then it is required to first save the permutation parameters. 
\\


\noindent \textbf{\underline{Results}}
\\

During `Model: Run' module, predictions, decision function values, weights and metrics measuring the model performance are computed. These results are stored in several subfields in \textcolor{blue}{PRT.output}.
\\

Within each main CV fold, statistics are computed at fold level. In versions 2.X, after looping over all folds, we concatenate predictions from each fold and pass them to \textcolor{red}{prt\_stats.m} together with true targets to compute model-level statistics. In v3.0, after looping over all folds, we compute the model level statistics by averaging fold-level stats across folds. In \textcolor{red}{prt\_stats.m}, for classification, we compute the confusion matrix, accuracy, balanced accuracy, accuracy by class, predictive value for each class, as well as AUC under ROC curve for binary classification. They are implemented by \textit{compute\_stats\_classifier} function. They call \textcolor{red}{prt\_tpr\_fpr.m} to compute AUC under ROC. For regression, we compute mean square error (MSE), normalized MSE, correlation between test and prediction, and squared correlation. They are implemented by \textit{compute\_stats\_regression} function. \textcolor{red}{prt\_stats.m} stores the values in structure stats and returns it to \textcolor{red}{prt\_cv\_model.m} and \textcolor{red}{prt\_nested\_cv.m}. 



%----------------------------------------------------------------------


\subsection{PRT fields created}

Two new fields are created in PRT.mat in this module: \textcolor{blue}{PRT.model.output.fold} and \textcolor{blue}{PRT.model.output.stats}. In \textcolor{blue}{PRT.model.output.fold}, several subfields are generated to store information of hyperparameter effects, prediction values in each fold, fold-level statistics, decision function values, model parameters and other information specific to some machines (e.g. kernel contributions for MKL models).
\\

%----------------------------------------------------------------------

\subsection{Functions called}

\textbf{\underline{Core:}}

\begin{enumerate}

\item \textcolor{red}{prt\_cv\_model.m}, \textcolor{red}{prt\_nested\_cv.m} and \textcolor{red}{prt\_cv\_fold.m}

\textcolor{red}{prt\_cv\_model.m} is where CV schemes are implemented within PRoNTo. It gets PRT.mat and user inputs from the interface.  Before looping over CV folds, it gets the model index, configure variables based on the CV matrix, gets values and variables for model estimation, such as user-specified targets, class numbers, kernels, and machine arguments. It also initializes outputs. \textcolor{red}{prt\_cv\_model.m} calls \textcolor{red}{prt\_getKernelModel} to loads kernels/features for kernel/non-kernel machines. \textcolor{red}{prt\_cv\_model.m} calls \textcolor{red}{prt\_nested\_cv.m} to optimize hyperparameters using inner CV by passing PRT.mat and the aforementioned fdata structure. \textcolor{red}{prt\_nested\_cv.m} returns the best hyperparameter value for external CV loops. \textcolor{red}{prt\_cv\_model.m} calls \textcolor{red}{prt\_cv\_fold.m} to train and test the predictive models by passing PRT.mat and fdata structure. \textcolor{red}{prt\_nested\_cv.m} also calls \textcolor{red}{prt\_cv\_fold.m} within each inner fold to estimate models. \textcolor{red}{prt\_cv\_fold.m} returns the model parameters and targets/labels. \textcolor{red}{prt\_cv\_model.m} also computes statistics for each fold and averaged statistics for the model by calling \textcolor{red}{prt\_stats.m}. 
       
\item \textcolor{red}{prt\_stats.m}

This function is called by \textcolor{red}{prt\_cv\_model.m} and \textcolor{red}{prt\_nested\_cv.m} to calculate fold-level and model-level metrics measuring the model performance. It takes predictions, model type, test targets, and class numbers as inputs.  For classification, the confusion matrix, accuracy, balanced accuracy, class accuracy, predictive value and AUC under ROC are calculated.  For regression, MSE (and normalized MSE), correlation and squared correlation are calculated. These metric values are passed as outputs in a ‘stats’ structure.
       
\item \textcolor{red}{prt\_machine.m} and \textcolor{red}{prt\_machine\_*.m} functions

These functions are in the \textit{machines} folder carrying out the actual modeling. \textcolor{red}{prt\_machine.m} is called by \textcolor{red}{prt\_cv\_fold.m} and passes data information and machine arguments to the machine-specific functions \textcolor{red}{prt\_machine\_*.m}, which are generally wrapper functions of the actual machines provided by several third-party libraries. These wrapper functions perform some initial checks on the inputs, re-structure data and arguments to formats accepted by these machines, then checks outputs from the estimated models and organize them into PRoNTo accepted structures. If developers want to add your own machines, make sure the input and output are in formats compatible to both your own machine and PRoNTo.  

\item \textcolor{red}{prt\_permutation.m}

This function carries out the permutation tests. It accepts user inputs from GUI or Batch interfaces, configures variables (e.g., CV matrix, fold numbers, and ID matrix), loads kernels/features for kernel/non-kernel machines, permutes data according to user inputs in Data \& Design module, runs models on the permuted data, calculates statistics for model performance, computes p-value  and saves the results.


\end{enumerate}


\noindent \textbf{\underline{GUI:}}

\begin{enumerate}
\setcounter{enumi}{4}

% \item \textcolor{red}{prt\_ui\_model.m} calls \textcolor{red}{prt\_cv\_model.m} to start the `Model: Run' module if the user clicks `Specify and Run'.

\item \textcolor{red}{prt\_ui\_cv\_model.m} and \textcolor{red}{prt\_ui\_cv\_model.fig}

This function pair define and implement the GUI interface for `Model: Run'. It loads the specified models from the PRT and fills a list that the user can select from. It calls \textcolor{red}{prt\_cv\_model.m} and \textcolor{red}{prt\_permutation.m} to execute the `Model: Run' module. \textcolor{red}{prt\_cv\_model} is called sequentially on the models selected for estimation.

\end{enumerate}

\noindent \textbf{\underline{Batch:}}

\begin{enumerate}
\setcounter{enumi}{5}

\item \textcolor{red}{prt\_cfg\_cv\_model.m} and \textcolor{red}{prt\_run\_cv\_model.m}

The configuration file defines the tree structure of `Model: Run' in Batch. It provides  interfaces for users to load PRT.mat, fills in permutation fields for users to choose, and provides fields for users to specify related parameters. It passes user inputs to \textcolor{red}{prt\_run\_cv\_model.m} for execution. \textcolor{red}{prt\_run\_cv\_model.m} loads PRT.mat and configure the specified model according to user inputs, then calls \textcolor{red}{prt\_cv\_model.m} to run model estimation. If permutation is selected, it calls \textcolor{red}{prt\_permutation.m} to perform the permutation tests.

\end{enumerate}


%======================================================================



%======================================================================

\section{Compute weights}


Weights are computed in two steps. They are carried out by \textcolor{red}{prt\_compute\_weights.m}, and \textcolor{red}{prt\_weights\_*.m} for specific machines. For NIfTI data, weight maps can be built as voxel-wise images or ROI-wise images. For mat and MEEG data types, weights are built in similar manners to those for NIfTI, though some specific changes have been made in related functions for the two types. Hence this section focus on the weight maps of NIfTI images in the following descriptions. For multiple kernel learning, maps indicating kernel contributions with ranking information can be obtained.
\\

\noindent \underline{\textbf{Weight computation}}
\\
 
The function \textcolor{red}{prt\_compute\_weights.m} accepts information from PRT.mat and some user-defined information, such as which model the weight image is for, where to save, flags indicating whether to build weights for each permutation and/or for each ROI, and atlas for ROI-wise weight image if it is selected. If different modalities were used without concatenating the samples, one weight map is built for each modality.
\\

For each modality, weight computation is classified into classification and regression problems. One weight image is built for each class in case of multi-class problems, whilst only one weight map is built for binary classification and regression problems. These are implemented in the scripts \textcolor{red}{prt\_compute\_weights\_class.m} and \textcolor{red}{prt\_compute\_weights\_regre.m}, respectively. Processing routines for classification and regression are similar to each other. Take classification as an example to illustrate the routine. \textcolor{red}{prt\_compute\_weights\_class.m} finds the type of data (NIfTI, .mat or MEEG) and machine first, then creates image names. One part needs some care is the multiple indexing step. In the `Data \& Design' module, users specify a 1st-level mask. The indices of features (e.g., voxels in NIfTI files) based on this mask are stored in \textcolor{blue}{PRT.fas} structure. In the feature set preparation step, if users have specified a second-level mask and/or atlas, indices of features used to build feature sets/kernels based on these are stored in PRT.fs structure. For example, in the `Prepare feature set' module, we need a subset of images to build the file array, from which we use potentially a subset to build a kernel. Furthermore, to avoid overloading the memory, the weight image is built slice by slice along the z axis, since most images are 3-D files. So we need to take indices from that slice, within both the second-level and the first-level masks. The three levels of indices are necessary to build weight maps. Hence, there is a need to perform the multi-level indexing search. This includes getting the correct indices from the 1st-level file array (.dat file) linking to each modality, from the 2nd-level feature sets (\textit{fs} structure), and from each slice along z axis, potentially within ROIs defined by 2nd-level mask and/or atlas.
\\
 
The next step is to create maps and check if weight images need to be created for the permutation test. Image initialization is performed before the weight computation. Weight images are created along z axis. One image is built for each fold, which results in a 4-D image where the last map is an average over all folds. Data operations are applied before weight computation. Then the data and information related to the machine used to generate the model are sent to \textcolor{red}{prt\_weights.m}. The code performing actual multiplication for weights is done by \textcolor{red}{prt\_weights.m} in the \textit{machines} folder within \textcolor{red}{prt\_compute\_weights\_class.m} and \textcolor{red}{prt\_compute\_weights\_regre.m}. This function calls other \textcolor{red}{prt\_weights\_*.m} in the \textit{machines} folder. Since different machines have different formulations, computing weights using outputs from these machines may vary too. Take \textcolor{red}{prt\_weights\_sMKL\_cla.m} as an example. The inputs are structure d and args similar to inputs to the machines. But the data loaded in \textcolor{red}{prt\_compute\_weights\_class.m} and \textcolor{red}{prt\_compute\_weights\_regre.m} is d.datamat. By taking coefficients from structure d and data from d.datamat, multiplications for weight values are carried out in these \textcolor{red}{prt\_weights\_*.m} functions. \textcolor{red}{prt\_weights.m} returns a weight vector to \textcolor{red}{prt\_compute\_weights\_class.m}.  Following procedures are to build the weights, reshape them into the correct dimensions, normalize the weights for visualization and save the images. In \textcolor{red}{prt\_compute\_weights.m}, the variable flag2 tells if we want to get one weight per region instead of one weight per voxel, so weight computation will be changed according to this. 


%----------------------------------------------------------------------


\subsection{PRT fields created}

In PRT.mat, extra fields are created: \textcolor{blue}{PRT.model.output.weight\_img}, \textcolor{blue}{PRT.model.output\_weight\_ROI} if building weights for ROIs is selected, subfields saving ranking information of weights, and \textcolor{blue}{PRT.model.output\_weight\_MOD} if weights for MKL are selected. 


%----------------------------------------------------------------------


\subsection{Files created}

This module creates weight image files: weight maps with `weights\_*' and/or `ROI\_weights\_*' in their names. The default naming rule is to use `weights\_' appended with the model name. If there are multiple classes or modalities, the class number or modality name will be appended to that name. Maps of ROIs are separated from voxel-wise maps which can be identified by names. If building weight images for permutations is selected, extra folders (with names perm\_* for voxelwise weights and perm\_ROI\_* for ROI-wise weights, respectively) containing weight maps per permutation are created.

% In PRT.mat, extra fields are created:

% \begin{itemize}
% \item \textcolor{blue}{PRT.model.output.weight\_img}, \textcolor{blue}{PRT.model.output\_weight\_ROI} if building weights for ROIs is selected and subfields saving information of rankings of weights.

% \item \textcolor{blue}{prt.model.output\_weight\_MOD} if weights for MKL are selected.

% \end{itemize}

% Display weights functionality searches for \textcolor{blue}{PRT.model.output.weight\_img} to locate where the maps are.


%----------------------------------------------------------------------


\subsection{Functions called}

% \textbf{\underline{Core:}}

\begin{enumerate}

\item \textcolor{red}{prt\_compute\_weights.m}, \textcolor{red}{prt\_compute\_weights\_class.m}, and \textcolor{red}{prt\_compute\_weights\_regre.m}
 
PRoNTo calls \textcolor{red}{prt\_compute\_weights.m} to implement weight image computation. It accepts several inputs: PRT.mat, user defined inputs including model name, image name, path to save the maps, atlas, one flag indicating whether to build weight images for permutations, and another flag indicating whether to build ROI-wise weight images. It outputs image files with specific names. This function computes the number of images needs to be built according to the number of modalities, and whether ROI-wise summary weight maps are selected by users. It finds the right model, loops over feature sets and/or modalities (over feature sets in v3.0, over modalities in versions 2.X, if building one kernel per modality) and saves PRT.mat. Within feature set loops, it further assesses the number of images according to whether it is a classification problem (class numbers) or regression problem. The actual weight map construction is carried out by \textcolor{red}{prt\_compute\_weights\_class.m} and  \textcolor{red}{prt\_compute\_weights\_regre.m} for classification and regression problems, respectively.  Computing weights for each voxel for a whole brain and for ROIs are processed differently by calling the two functions. If ROI-wise weight images are selected to be built, a flag `flag2' will be set to 1 and sent to the two functions. 

\bi
\item \textcolor{red}{prt\_compute\_weights\_class.m} and \textcolor{red}{prt\_compute\_weights\_regre.m} also deal with weight computations for permutations. The following procedures are initializing the image, building images slice by slice along the z axis, building one image per fold, applying operations, building weights (a 4D image, the last dimension is the number of folds plus one average map across folds), normalizing the weights for visualization and save the images. They call \textcolor{red}{prt\_weights.m} and \textcolor{red}{prt\_weights\_*.m} to perform the actual multiplications for weight computations.  

\item \textcolor{red}{prt\_compute\_weights\_class.m} and \textcolor{red}{prt\_compute\_weights\_regre.m} save weight maps and return the image names to \textcolor{red}{prt\_compute\_weights.m}. \textcolor{red}{prt\_compute\_weights.m} fills weight fields in \textcolor{blue}{PRT.model.output} and updates PRT.mat. 
\ei 

\item \textcolor{red}{prt\_weights.m}, \textcolor{red}{prt\_weights\_bin\_linkernel.m}, and \textcolor{red}{prt\_weights\_*.m} for other machines

These functions are in the machines folder. Formulations of different machines may vary, so weight computations may be machine-specific. For the kernel machines, the coefficients are obtained from the kernel space, then mapped back to the original space. When adding new machines, corresponding weight computations shall be considered and added too. There are some general cases in PRoNTo, where \textcolor{red}{prt\_weights\_bin\_linkernel.m} can be used. Other \textcolor{red}{prt\_weights\_*.m} functions implement multiplications for machines beyond the linear kernel case.
\\
 
One example of the machine-specific weight function is \textcolor{red}{prt\_weights\_sMKL\_cla.m}. It accepts input structures d and args, creates weight vector as the output. A list of subfields in structure d is shown as follows:

\begin{itemize}
\item d.datamat: the data being loaded in \textcolor{red}{prt\_compute\_weights\_class.m} and \textcolor{red}{prt\_compute\_weights\_regre.m} of dimension number of $\#$features x $\#$number of examples, corresponding to the features in the training set.
\item d.coeffs:– coefficients, such as alphas generated by using SVM.
\item d.betas: this is the kernel contributions, specific to sMKL.
\item d.idfeat\_img: voxel indices of the features in the image.
\end{itemize}

\end{enumerate}



%======================================================================



%======================================================================

% \section{Results}

% During `Model: Run' module, predictions, decision function values, weights and metrics measuring the model performance are computed. These results are stored in several subfields in \textcolor{blue}{PRT.output}.
% \\

% \underline{\textbf{How are the performance metrics computed?}}
% \\

% Within each main CV fold, statistics are computed at fold level. After looping over all folds, we concatenate predictions from each fold and pass them to \textcolor{red}{prt\_stats.m} together with true targets to compute model-level statistics. Formulations and codes calculating these metrics are in \textcolor{red}{prt\_stats.m}. ROC and AUC for binary classification are calculated in Display results module by calling \textcolor{red}{prt\_plot\_ROC.m} in versions 2.X. This is moved to \textcolor{red}{prt\_stats.m} from v3.0.


% %----------------------------------------------------------------------

% \subsection{Functions called}

% % \textbf{\underline{Core:}}

% \begin{enumerate}

% \item \textcolor{red}{prt\_stats.m}

% The same as described in the `Compute weights' module. 

% \item \textcolor{red}{prt\_plot\_ROC.m}

% This function takes PRT, model index, fold and graphical information as inputs, then obtains true positive and false positive rates to build ROC curves and calculates AUC. This is done for each fold. Averaged results over folds are given too.  

% \end{enumerate}





